本付録では、本実験4（時間重み自動判別）に関して、本文から切り出した詳細結果を示す。

\subsubsection*{時間パターン別の優位性（18名データ）}
時間パターンの優位性をより俯瞰するため、\textbf{学習ステップ20（履歴20件）}において、全重み付けパターン（6種）における各時間パターンの平均Top-Kを集計した。その結果、\textbf{直近7日指数減衰}や\textbf{Softmax 7日}、\textbf{セッション線形減衰}が上位に並び、短期の情報を重視する設計が安定して高い精度を示した（表\ref{tab:time-pattern-ranking-20}）。これは本実験3の結論（古い履歴の混入が精度低下に繋がる）と整合的である。ただし、上位方式間の差は僅差であり、ワーカーごとの最適が分散する可能性も高い。このため、自動選択は「短期方式を常に採用する」ことではなく、\textbf{必要なときに短期へ寄せる}ためのスイッチとして設計することが望ましい。

\begin{table}[htbp]
\centering
\caption{18名データの時間パターン上位（学習20件, 平均Top-K, 全重み付け平均）}
\label{tab:time-pattern-ranking-20}
\scriptsize
\begin{tabular}{l l r r r}
\hline
順位 & 時間パターン & 平均Top-1 & 平均Top-3 & 平均Top-5 \\
\hline
1 & 直近7日指数減衰 & 0.5352 & 0.6880 & 0.7296 \\
2 & Softmax 7日 & 0.5269 & 0.6778 & 0.7269 \\
3 & セッション順線形10件 & 0.5278 & 0.6731 & 0.7250 \\
4 & セッション順線形20件 & 0.5139 & 0.6324 & 0.7241 \\
5 & 直近30日線形減衰 & 0.5278 & 0.6824 & 0.7222 \\
\hline
\end{tabular}
\end{table}

また、\texttt{equal} のみで各ワーカーの最良パターンを確認すると、時間重み自動判別が11名、直近7日指数減衰が3名で最多となり、残りは他の減衰方式に分散した。これは\textbf{集団平均では短期重視が有利である一方、個別最適は多様}であることを示している。

\subsubsection*{20名データでの追試（drift 閾値の影響）}
本実験4の実装は、drift 閾値の設定により「安定（長期）」と「変化（短期）」の切替が変わる。そこで、20名データ（うち1名はデータ読込エラーのため除外し、19名で評価）に対して、drift 閾値を変えた複数条件で実行し、学習ステップ別の平均Top-Kを比較した。なお、学習ステップ25以降は履歴が不足するワーカーが増えるため、評価対象人数が減少する。

\paragraph{drift 閾値 0.5 / 0.7}
結果として、両設定で精度はほぼ一致し、時間重み自動判別は\textbf{全ワーカーが安定（長期）判定}となり、採用方式は\textbf{期間制限型（直近90日ウィンドウ）}に揃った。drift の分布は最小 0.005、最大 0.154（中央値 0.035）であり、いずれの閾値でも短期側へ切り替わらない範囲に収まっていた。学習ステップ別の平均Top-Kを表\ref{tab:time-auto-20workers-steps}に示す。

\begin{table}[htbp]
\centering
\caption{20名データにおける時間重み自動判別（\texttt{equal}）の学習ステップ別平均Top-K（drift 閾値 0.5/0.7）}
\label{tab:time-auto-20workers-steps}
\scriptsize
\begin{tabular}{c r r r r}
\hline
学習ステップ & 対象人数 & 平均Top-1 & 平均Top-3 & 平均Top-5 \\
\hline
5  & 19 & 0.5316 & 0.6158 & 0.6947 \\
10 & 19 & 0.4526 & 0.5526 & 0.6211 \\
15 & 19 & 0.5105 & 0.6053 & 0.6895 \\
20 & 19 & 0.4789 & 0.6316 & 0.7000 \\
25 & 15 & 0.4733 & 0.5800 & 0.6333 \\
30 & 11 & 0.3545 & 0.4636 & 0.5273 \\
35 & 11 & 0.3455 & 0.5182 & 0.5455 \\
40 & 8  & 0.4125 & 0.4750 & 0.6500 \\
\hline
\end{tabular}
\end{table}

\paragraph{drift 閾値 0.05}
drift 閾値を 0.05 に設定すると、19名のうち\textbf{安定（長期）が11名、直近（短期）が8名}となり、採用方式は\textbf{期間制限型（直近90日ウィンドウ）}と\textbf{直近7日指数減衰}に分かれた。学習ステップ別の平均Top-Kを表\ref{tab:time-auto-20workers-steps-thr005}に示す。

\begin{table}[htbp]
\centering
\caption{20名データにおける時間重み自動判別（\texttt{equal}）の学習ステップ別平均Top-K（drift 閾値 0.05）}
\label{tab:time-auto-20workers-steps-thr005}
\scriptsize
\begin{tabular}{c r r r r}
\hline
学習ステップ & 対象人数 & 平均Top-1 & 平均Top-3 & 平均Top-5 \\
\hline
5  & 19 & 0.5316 & 0.6158 & 0.6947 \\
10 & 19 & 0.4632 & 0.5474 & 0.6158 \\
15 & 19 & 0.5053 & 0.6158 & 0.7000 \\
20 & 19 & 0.4789 & 0.6316 & 0.7000 \\
25 & 15 & 0.4733 & 0.5800 & 0.6333 \\
30 & 11 & 0.3545 & 0.4636 & 0.5273 \\
35 & 11 & 0.3455 & 0.5182 & 0.5455 \\
40 & 8  & 0.3250 & 0.4750 & 0.6500 \\
\hline
\end{tabular}
\end{table}

\paragraph{drift 閾値 0.03}
さらに drift 閾値を 0.03 に設定すると、\textbf{直近（短期）が11名、安定（長期）が8名}となり、短期側へ切り替わるワーカーが増加した。学習ステップ別の平均Top-Kを表\ref{tab:time-auto-20workers-steps-thr003}に示す。

\begin{table}[htbp]
\centering
\caption{20名データにおける時間重み自動判別（\texttt{equal}）の学習ステップ別平均Top-K（drift 閾値 0.03）}
\label{tab:time-auto-20workers-steps-thr003}
\scriptsize
\begin{tabular}{c r r r r}
\hline
学習ステップ & 対象人数 & 平均Top-1 & 平均Top-3 & 平均Top-5 \\
\hline
5  & 19 & 0.5316 & 0.6158 & 0.6947 \\
10 & 19 & 0.4632 & 0.5421 & 0.6263 \\
15 & 19 & 0.5000 & 0.6526 & 0.7158 \\
20 & 19 & 0.5211 & 0.6632 & 0.7316 \\
25 & 15 & 0.4600 & 0.5800 & 0.6667 \\
30 & 11 & 0.3545 & 0.4636 & 0.5273 \\
35 & 11 & 0.3455 & 0.5000 & 0.5455 \\
40 & 8  & 0.3250 & 0.4500 & 0.6375 \\
\hline
\end{tabular}
\end{table}
